\documentclass[10pt,a4paper]{amsart}
\usepackage{amsmath,amssymb,amsthm,graphicx}
% Retained minimal preamble for the Stacks Project chapter topology.tex (Topology).
% The Stacks Project uses one single shared preamble for every chapter; the snapshot stages it
% identically as sources/stacks-topology/preamble.tex, sources/stacks-chow/preamble.tex and
% sources/stacks-etale-cohomology/preamble.tex.  The block below is copied mechanically, line for
% line, from that shared preamble: the Xy-pic setup (lines 14-18) and the native theorem
% declarations together with the mathematical macros (lines 154-217), plus the project's own
% reference/slogan/history display environments (lines 8-12, via the verbatim package).
% Nothing was added, removed or redefined; the wrapper class and the amsmath/amssymb/amsthm
% base packages are supplied by the assembler.
% Chapter: topology.tex -- Topology.
\usepackage[all]{xy}

% We use 2cell for 2-commutative diagrams.
\xyoption{2cell}

% Theorem environments.
%
\theoremstyle{plain}
\newtheorem{theorem}[subsection]{Theorem}
\newtheorem{proposition}[subsection]{Proposition}
\newtheorem{lemma}[subsection]{Lemma}

\theoremstyle{definition}
\newtheorem{definition}[subsection]{Definition}
\newtheorem{example}[subsection]{Example}
\newtheorem{exercise}[subsection]{Exercise}
\newtheorem{situation}[subsection]{Situation}

\theoremstyle{remark}
\newtheorem{remark}[subsection]{Remark}
\newtheorem{remarks}[subsection]{Remarks}

\numberwithin{equation}{subsection}

% Macros
%
\def\lim{\mathop{\mathrm{lim}}\nolimits}
\def\colim{\mathop{\mathrm{colim}}\nolimits}
\def\Spec{\mathop{\mathrm{Spec}}}
\def\Hom{\mathop{\mathrm{Hom}}\nolimits}
\def\Ext{\mathop{\mathrm{Ext}}\nolimits}
\def\SheafHom{\mathop{\mathcal{H}\!\mathit{om}}\nolimits}
\def\SheafExt{\mathop{\mathcal{E}\!\mathit{xt}}\nolimits}
\def\Sch{\mathit{Sch}}
\def\Mor{\mathop{\mathrm{Mor}}\nolimits}
\def\Ob{\mathop{\mathrm{Ob}}\nolimits}
\def\Sh{\mathop{\mathit{Sh}}\nolimits}
\def\NL{\mathop{N\!L}\nolimits}
\def\CH{\mathop{\mathrm{CH}}\nolimits}
\def\proetale{{pro\text{-}\acute{e}tale}}
\def\etale{{\acute{e}tale}}
\def\QCoh{\mathit{QCoh}}
\def\Ker{\mathop{\mathrm{Ker}}}
\def\Im{\mathop{\mathrm{Im}}}
\def\Coker{\mathop{\mathrm{Coker}}}
\def\Coim{\mathop{\mathrm{Coim}}}

% Boxtimes
%
\DeclareMathSymbol{\boxtimes}{\mathbin}{AMSa}{"02}

%
% Macros for moduli stacks/spaces
%
\def\QCohstack{\mathcal{QC}\!\mathit{oh}}
\def\Cohstack{\mathcal{C}\!\mathit{oh}}
\def\Spacesstack{\mathcal{S}\!\mathit{paces}}
\def\Quotfunctor{\mathrm{Quot}}
\def\Hilbfunctor{\mathrm{Hilb}}
\def\Curvesstack{\mathcal{C}\!\mathit{urves}}
\def\Polarizedstack{\mathcal{P}\!\mathit{olarized}}
\def\Complexesstack{\mathcal{C}\!\mathit{omplexes}}
% \Pic is the operator that assigns to X its picard group, usage \Pic(X)
% \Picardstack_{X/B} denotes the Picard stack of X over B
% \Picardfunctor_{X/B} denotes the Picard functor of X over B
\def\Pic{\mathop{\mathrm{Pic}}\nolimits}
\def\Picardstack{\mathcal{P}\!\mathit{ic}}
\def\Picardfunctor{\mathrm{Pic}}
\def\Deformationcategory{\mathcal{D}\!\mathit{ef}}

% --- appended 2026-10-12 (Chief): the project own definitions of the reference/slogan/history
% display environments, copied mechanically from the shared Stacks preamble lines 8-12
% (\usepackage{verbatim} plus the three \newenvironment lines, which the chapter preamble keeps
% verbatim from the project preamble). Nothing added, removed or redefined: these environments
% carry no mathematical content and the project itself typesets them as comments. Without these
% three lines a counted statement whose own text uses them fails with Environment slogan undefined.
\usepackage{verbatim}
\newenvironment{reference}{\comment}{\endcomment}
\newenvironment{slogan}{\comment}{\endcomment}
\newenvironment{history}{\comment}{\endcomment}
\title{Opens and quasi-compact opens of a cofiltered limit of topological spaces}
\author{The Stacks Project Authors}
\date{stacks-project:topology.tex:local-snapshot-20261012; CC BY 4.0}
\begin{document}
\maketitle
\noindent\textit{Editorial note.} All delivered mathematical text is the author's own, copied from the pinned Stacks Project chapter source sources/stacks-topology/pinned-original-source.tex; no statement or proof step was rewritten, reordered or removed.\par\medskip
\noindent\textit{Editorial note.} Preamble provenance: the Stacks Project keeps one single shared preamble for every chapter, staged here as sources/stacks-topology/preamble.tex; this card uses retained-preamble.tex, which copies that shared preamble mechanically (the Xy-pic setup, the native theorem declarations, the mathematical macros and the reference/slogan/history display environments). Nothing was added, removed or redefined; the wrapper class and the base packages are supplied by the assembler.\par\medskip
\noindent\textit{Editorial note.} Bibliography provenance: the chapter carries no chapter-local bibliography (the project keeps one shared bibliography file that is not staged here) and neither the counted statement nor its proof contains any citation, so the card carries no bibliography and no entry was reconstructed.\par\medskip
\noindent\textit{Editorial note.} This card was selected from the staged candidate list for this chapter (candidates-halfA.json / candidates-halfB.json), where it is the candidate labelled lemma-describe-limits; the statement and proof spans are the candidate's own statement\_lines and proof\_lines, checked against the source so that each span opens and closes a real, uncommented environment.\par\medskip
\noindent\textit{Allowed standard prerequisites (editorial).}\begin{itemize}
\item The construction of the limit of a diagram of topological spaces over a cofiltered category, as the subset of the product with the induced topology. Given in the author's section on limits of spaces, immediately before the counted statement. Supplies the description of the limit and its projection maps in terms of which both assertions are formulated.
\end{itemize}\medskip
\begin{lemma}
\label{lemma-describe-limits}
Let $\mathcal{I}$ be a cofiltered category. Let $i \mapsto X_i$ be a diagram
of topological spaces over $\mathcal{I}$. Let $X = \lim X_i$ be the limit
with projection maps $f_i : X \to X_i$.
\begin{enumerate}
\item Any open of $X$ is of the form $\bigcup_{j \in J} f_j^{-1}(U_j)$
for some subset $J \subset I$ and opens $U_j \subset X_j$.
\item Any quasi-compact open of $X$ is of the form
$f_i^{-1}(U_i)$ for some $i$ and some $U_i \subset X_i$ open.
\end{enumerate}
\end{lemma}

\begin{proof}
The construction of the limit given above shows that $X \subset \prod X_i$
with the induced topology. A basis for the topology of $\prod X_i$ are
the opens $\prod U_i$ where $U_i \subset X_i$ is open and $U_i = X_i$
for almost all $i$. Say $i_1, \ldots, i_n \in \Ob(\mathcal{I})$ are the
objects such that $U_{i_j} \not = X_{i_j}$. Then
$$
X \cap \prod U_i = f_{i_1}^{-1}(U_{i_1}) \cap \ldots \cap f_{i_n}^{-1}(U_{i_n})
$$
For a general limit of topological spaces these form a basis for the
topology on $X$. However, if $\mathcal{I}$ is cofiltered as in the statement
of the lemma, then we can pick a $j \in \Ob(\mathcal{I})$ and morphisms
$j \to i_l$, $l = 1, \ldots, n$. Let
$$
U_j =
(X_j \to X_{i_1})^{-1}(U_{i_1}) \cap \ldots \cap
(X_j \to X_{i_n})^{-1}(U_{i_n})
$$
Then it is clear that $X \cap \prod U_i = f_j^{-1}(U_j)$. Thus for any open
$W$ of $X$ there is a set $A$ and a map $\alpha : A \to \Ob(\mathcal{I})$ and
opens $U_a \subset X_{\alpha(a)}$ such that
$W = \bigcup f_{\alpha(a)}^{-1}(U_a)$. Set $J = \Im(\alpha)$ and
for $j \in J$ set $U_j = \bigcup_{\alpha(a) = j} U_a$ to see that
$W = \bigcup_{j \in J} f_j^{-1}(U_j)$.
This proves (1).

\medskip\noindent
To see (2) suppose that $\bigcup_{j \in J} f_j^{-1}(U_j)$ is quasi-compact.
Then it is equal to
$f_{j_1}^{-1}(U_{j_1}) \cup \ldots \cup f_{j_m}^{-1}(U_{j_m})$
for some $j_1, \ldots, j_m \in J$. Since $\mathcal{I}$ is cofiltered,
we can pick a $i \in \Ob(\mathcal{I})$ and morphisms
$i \to j_l$, $l = 1, \ldots, m$. Let
$$
U_i =
(X_i \to X_{j_1})^{-1}(U_{j_1}) \cup \ldots \cup
(X_i \to X_{j_m})^{-1}(U_{j_m})
$$
Then our open equals $f_i^{-1}(U_i)$ as desired. 
\end{proof}

\end{document}
